\documentclass[12pt, letterpaper]{article}

\usepackage[utf8]{inputenc}
\title{Week 2}
\author{Raadhikka Gupta}
\date{\today}

\begin{document}
\maketitle

\section{Prompt Engineering}

\subsection{Type of prompts}
\begin{itemize}
    \item \textbf{Direct Prompts}: clear and specific instructions; best for simple tasks when you know exactly what output you need the AI to generate.
    \item \textbf{Open-ended Prompts}: diverse outputs without any constraints; for brainstorming, storytelling, or exploratory discussions.
    \item \textbf{Task-specific Prompts}: precise, goal-oriented tasks; defined focus and have additional context or examples to help generate responses.
\end{itemize}

\subsection{Prompt Engineering Techniques}
\begin{itemize}
    \item \textbf{Zero-shot prompting}: response based on pre-trained knowledge without any examples or hints in the prompt; for simple, familiar, and straightforward tasks. 
        \begin{itemize}
            \item They can be less effective for tasks with domain knowledge, high understanding and specific outputs; they can have biases in training that can affect fairness and accuracy.
        \end{itemize}
    \item \textbf{One-shot prompting}: model is given a single example along with a task prompt to guide its response; for cases where collecting large amounts of data is impractical. 
        \begin{itemize}
            \item Add sample input and output, and the instruction in the prompt.
            \item They can have biases because one example shape the output, and not work for complex tasks or tasks that need domain knowledge.
        \end{itemize}
    \item \textbf{Few-shot prompting}: model is given a few examples along with a task prompt to guide its performance; helps to find patterns and improve accuracy of response.
        \begin{itemize}
            \item Add sample inputs and outputs, and the instruction in the prompt.
            \item They need to designed carefully and require domain knowledge to create effective examples so that the model can pick up the right things; more intensive as more examples increase processing time.
        \end{itemize}
    \item \textbf{Chain-of-thought prompting}: technique to enhance the output of LLM by guiding a model through a step-by-step reasoning process by using a coherent series of logical steps; useful for when prompting the model to perform tasks with logic or maths.
        \begin{itemize}
            \item Use a checklist and include these in the prompt: i) a clear task or problem statement, ii) a reasoning trigger or phrase that cues AI to reason, iii) a step-by-step format to organize the reasoning, iv) a format to summarize the final answer.
            \item They are less effective for tasks with complex or branching logic as the model follows a liner reasoning chain; best with instructions are clear and guide effectively, as vague instructions = poor output.
        \end{itemize}
    \item \textbf{Meta prompting}: technique where model is asked to generate or improve its own prompt; to help enhance the quality, focus or consistency of output so that it will be more accurate, polished with the required tone, style or level of detail.
        \begin{itemize}
            \item Add instruction and guidance for refinement to the prompt.
            \item They need well-structured guidance to refine them; add complexity and processing time as the model should analyze and modify a prompt before generating a response.
        \end{itemize}
\end{itemize}

\subsection{Mistakes to avoid when prompting}
\begin{itemize}
    \item Vague Prompts
    \item Grammatical errors
    \item Leading questions
    \item Context gaps
    \item Complex Prompts
\end{itemize}

\subsection{Best Practices}
\begin{itemize}
    \item Start with a clear objective
    \item Provide context
    \item Be specific and precise
    \item Use a conversational tone
    \item Refine and iterate
\end{itemize}

\section{Chatbots}
\subsection{Types of Chatbots}
\begin{itemize}
    \item \textbf{Text-based}: queries only through text-based interactions; help in customer service handling various routine queries like order status, and detail updates.
    \item \textbf{Voice-based}: spoken commands using speech recognition to interpret spoken query and responds verbally; can handle complex interactions and adapt to various accents and speech patterns.
    \item \textbf{Virtual assistants}: provide a smooth and multifunctional UX; for a variety of tasks like answering emails, tracking orders, scheduling appointments, weather updates and controlling smart home devices.
\end{itemize}

\subsection{Conversational AI}
\begin{itemize}
    \item It enables machines to engage in human-like conversations with users.
    \item Uses NLP, machine learning, and sentiment analysis.
        \begin{itemize}
            \item \textbf{Natural language Processing}: understand and interpret human language; to process the words and phrases you use, understand the meaning, and determine the context.
            \item \textbf{Machine learning}: learn from data and improve its performance over time by learning from past interactions, recognize patterns, and improve responses over time.
            \item \textbf{Sentiment analysis}: detect and understand the emotional tone; mood detector and enables to identify whether the message is positive, negative, or neutral.
        \end{itemize}
    \item This allows the chatbot to perform tasks like interpreting human language, remembering past conversations, incorporating context into their responses, and creating more intuitive and effective interactions.
\end{itemize}

\subsection{Designing Chatbots}
\begin{itemize}
    \item \textbf{User Interface}: designing the chatbot personality, and engagement
        \begin{itemize}
            \item \textbf{Chatbot personality}: giving it unique characteristics and communication style that aligns with the company identity; good one helps users connect with the chatbot and improves their experience.
                \begin{itemize}
                    \item \textbf{Appearance}: layout such as colours, buttons, and fonts.
                    \item \textbf{Interaction}: user experience design; focusing on the user's interaction with the chatbot, the information that they have, and the information that they recieve.
                    \item \textbf{Dialog}: responses for the chatbot; writing clear, helpful, and engaging text and matching the company brand with the tone and personality.
                    \item \textbf{Learning}: the ability to adapt to different user actions; involves machine learning which helps the chatbot to interpret accurately and give personalized responses; with NLP and ML, you train the chatbot to interpret users' question and answers and respond.
                \end{itemize}
            \item \textbf{Engagement}: making the chatbot interactive and easy to use so that uses want to interact with it; easy to use, responsive, interactive, and reflection of the company brand.
        \end{itemize}
    \item \textbf{Meeting customer needs}
        \begin{itemize}
            \item \textbf{Define the goal}: what you want it to do; keeps the focus on a specific need 
            \item \textbf{Choose the starting domain}: specific area to focus on initially; eg billing section of a website, so domain to integrate the chatbot to effectively.
            \item \textbf{Choose the starting channel}: platform that is the most frequently visited by customers.
            \item \textbf{Start small}: assigning a small, manageable task to the bot to test its functionality; how well it performs and gather feedback.
            \item \textbf{Identify key topics}: focus on few most-asked specific questions or tasks that it will handle to ensure that it give the most accurate information for them.
            \item \textbf{Identify relevant content sources}: take stock of all the helpful content such as product information, articles, and FAQs within the domain, and set up a quick search function.
            \item \textbf{Plan for handoff}: when chatbot can not handle a request, plan to escalate to a human agent.
        \end{itemize}
    \item \textbf{Design flow logic}: how it reacts to each statement made during an exchange and proceeds to the next step; flow logic needs simple if-else statements, decision trees, complex algorithms, or machine learning-driven probabilistic logic.
\end{itemize}

\end{document}